\section{The \CEGMem\ Architecture}
\label{sec:arch}

\subsection{Problem setup}
A repair task is a tuple $(\cand, \spec, \Ok)$: a candidate space of patches, a
specification, and an oracle. A patch is \emph{semantically correct} iff it
satisfies $\spec$ on every input in a domain $\mathcal{X}$, and we write
$\Gstar$ for that set, assumed non-empty. No implementable oracle decides
membership in $\Gstar$. What we have is a \emph{counterexample oracle at depth
$\oracledepth$}: a call draws a check set $\Xk$ of
$\min(\oracledepth, \lvert\mathrm{pool}\rvert)$ inputs from the fault's shipped
pool and runs them until one separates candidate from reference,
\begin{equation}
\Ok(p) =
\begin{cases}
x \in \Xk \text{ with } \lnot\spec(p, x) & \text{if a checked input fails},\\
\textsf{accept} & \text{otherwise}.
\end{cases}
\label{eq:oracle}
\end{equation}
Write $\Gk$ for what one such call certifies---the patches passing every
$x \in \Xk$---and $\Gpool$ for those passing the \emph{entire} pool. Two
properties of $\Xk$ are load-bearing. It is drawn \emph{per call}, so $\Gk$ is
a family indexed by the draw; but the draw is degenerate whenever
$\oracledepth \ge \lvert\mathrm{pool}\rvert$, where $\Xk$ is the pool and
$\Gk = \Gpool$---the case for $97$ of our $99$ faults at the
$\oracledepth = 100$ every headline number uses. The three
sets nest, $\Gstar \subseteq \Gpool \subseteq \Gk$, so acceptance is the
weakest of them and $\Gk \setminus \Gstar$ is \emph{false acceptance};
\cref{sec:results-depth} measures that gap and never assumes it away. Three budget quantities also stay distinct: the proposal budget
$\budget$, the oracle budget $b$ against which repair rate is reported, and the
oracle depth $\oracledepth$.

\ifsuppmatter\else\begin{table}[t]
\centering
\footnotesize
\caption{Notation. The three horizontal blocks are the distinctions this paper
depends on and that its earlier drafts blurred: what is correct versus what an
oracle certifies; what is computable before an execution versus after one; and
which budget a number is denominated in.}
\label{tab:notation}
\setlength{\tabcolsep}{4pt}
\begin{tabular}{@{}ll@{}}
\toprule
\textbf{Symbol} & \textbf{Meaning} \\
\midrule
$\cand, p$ & candidate space; a candidate patch \\
$\spec, \mathrm{sat}(p)$ & specification; $p$ satisfies $\spec$ on all of $\mathcal{X}$ \\
$\prop$ & stochastic proposer (the LLM) \\
$\pi,\ q_\tau$ & base mass on the correct class; on type $\tau$ \\
$\wmass_s$ & base mass of class $s$ ($\wmass_\top{=}\pi$, $\wmass_\tau{=}q_\tau$) \\
\midrule
$\Gstar \subseteq \cand$ & \emph{semantically} correct patches \\
$\Gpool$ & patches passing the entire shipped test pool \\
$\Gk$ & patches the depth-$\oracledepth$ oracle accepts;
  $\Gstar {\subseteq} \Gpool {\subseteq} \Gk$ \\
$\Ok(p)$ & counterexample oracle at depth $\oracledepth$ (\cref{eq:oracle}) \\
\midrule
$\loc(p)$ & edit location: the key computable \emph{before} execution \\
$\ttype(p, x) \in \types$ & failure label, defined \emph{after} $x$ refutes $p$ \\
$\mem,\ \mem[\ell]$ & memory of triples $(p, x, \tau)$; its $\loc{=}\ell$ bucket \\
$\elim \subseteq \types$ & labels this arm has eliminated \\
$\elimobs \subseteq \types$ & labels refuted so far, arm-neutral \\
$\Kreach$ & reachable failure types, $|\{\tau : q_\tau > 0\}|$ \\
$\info,\ \coh$ & oracle informativeness; typing coherence (\cref{def:coh}) \\
\midrule
$\budget$ & \emph{proposal} budget per episode ($20$ throughout) \\
$b$ & \emph{oracle-call} budget: the reported cost axis \\
$\oracledepth$ & \emph{oracle depth}: test cases per verification \\
\bottomrule
\end{tabular}
\end{table}
\fi

\subsection{Failure labels, and the key available before execution}
A refuted candidate $p$ with counterexample $x$ receives a label
\begin{equation}
\ttype(p, x) = \big(\loc(p),\ \text{property violated at } x\big) \in \types ,
\label{eq:type}
\end{equation}
where $\loc(p)$ is the patch's \emph{edit location}, obtained by diffing the
candidate against the buggy source. The two components differ in when they are
available, and that is load-bearing: $\loc(p)$ needs no execution, while the
violated property is only defined \emph{after} a failing run and so does not
exist for a fresh candidate. The guard therefore cannot index on $\ttype$. It
indexes on $\loc$, and memory's buckets are keyed on it:
$\mem[\ell] = \{(p', x', \tau') \in \mem : \loc(p') = \ell\}$. The full label
$\ttype$ is what steering excludes---it names classes already \emph{observed},
for which the property component exists---and what the coherence sweep
perturbs. Nothing here requires predicting a violated property before running
anything.

\begin{definition}[Informativeness and coherence]
\label{def:coh}
The oracle is $\info$-\emph{informative} if a counterexample refuting a
candidate of class $\tau$ refutes every candidate of class $\tau$ with
probability at least $\info$. Memory is $\coh$-\emph{coherent} if it attributes
a candidate or counterexample to its true class with probability at least
$\coh$. The ideal regime is $\info = \coh = 1$.
\end{definition}

\subsection{Typed memory: guard and steer}
\CEGMem\ maintains a memory $\mem$ of triples $(p, x, \tau)$ and the induced
set of eliminated classes $\elim$. It exposes two operations, kept distinct
because \cref{sec:results} shows they do different jobs and that only one
works.

\paragraph{Guard.}
Before paying the oracle, the guard rejects a candidate that a stored
counterexample still refutes:
\begin{equation}
\mathrm{Guard}_\mem(p) = \mathbb{1}\!\left[\exists (p', x', \tau') \in \mem :
x' \text{ refutes } p\right].
\label{eq:guard}
\end{equation}
Here $x'$ \emph{refutes} $p$ under the oracle's own predicate---$p$ on $x'$
disagrees with the reference, crashes, or times out---so a block and a
refutation are the same event, evaluated by different callers.
An \emph{indexed} guard evaluates \cref{eq:guard} against $\mem[\loc(p)]$ first
and then the rest of memory---the fallback belongs to the analyzed algorithm
and to the shipped one, not to one of them---while an \emph{untyped} guard
walks memory in insertion order. Both short-circuit at the first entry that still refutes and both
consult all of memory when none does, so their worst cases coincide: the index
changes the \emph{order} of consultation, not its bound (\Rguardcost). Our
implementation additionally deduplicates stored counterexamples by input and
memoizes verdicts within a round, reducing executions at unchanged
consultations.

Two properties of \cref{eq:guard} are easy to lose. It is a \emph{verified}
predicate, not a match heuristic: a block happens only when a stored
counterexample is re-executed and still fails, so widening the search can never
cost soundness (\cref{thm:sound}). And the index must say where to look
\emph{first}, not where to \emph{stop}: $\loc$ records where a candidate edits
while refutation is decided by which input it gets wrong, so a bucket-only
guard is an unsound stopping rule. \Cref{sec:impl} reports what that cost us.

\paragraph{Steer.}
Memory also reshapes the proposer. Let $\wmass_s$ be its base mass on class
$s$, with $\wmass_\top = \pi$ for the correct class and $\wmass_\tau = q_\tau$
otherwise. The steered proposer renormalizes onto the surviving support:
\begin{equation}
\mathbb{P}_{\mathrm{steer}}[s] =
\frac{\wmass_s \cdot \mathbb{1}[s \notin \elim]}
     {\pi + \sum_{\tau \notin \elim} q_\tau},
\qquad s \in \{\top\} \cup \types .
\label{eq:steer}
\end{equation}
This is what a flat transcript cannot do---excluding a class requires having
one---and it idealizes an operation that, with an LLM, can only be
\emph{requested}: an exclusion instruction plus class-matched evidence in the
prompt. Whether the request is honored is empirical;
\cref{sec:results-ablation} answers no for a 7B proposer, while the same
renormalization applied on the guard side, by not charging blocked rounds, does
deliver (\cref{sec:results-e9}).

\subsection{The repair loop}
Each round draws a steered proposal, guards it against memory, and calls the
oracle only if it survives; an accepted patch is returned and a refuted one
stored with its counterexample and label, eliminating that class. The arms in
\cref{sec:results} switch the two operations independently
(\appref{alg:cegmem}{the supplement} gives the loop in full). \CEGMem\ adds no
planner and no controller---the only new state is the memory---which is what
keeps \cref{sec:theory}'s guarantees properties of \emph{memory}, and so
independent of the proposer.

\ifsuppmatter\else\begin{algorithm}[t]
\caption{\CEGMem: counterexample-guided repair with typed memory. The two
memory operations are lines~\ref{alg:steer} (steer) and \ref{alg:guard}
(guard); switching each off independently yields the five arms of
\cref{sec:design}.}
\label{alg:cegmem}
\DontPrintSemicolon
\SetKwInOut{Input}{Input}
\SetKwInOut{Output}{Output}
\Input{task $(\cand, \spec, \Ok)$, proposer $\prop$, budget $\budget$}
\Output{a repaired patch, or $\bot$}
$\mem \leftarrow \emptyset$; $\elim \leftarrow \emptyset$
  \tcp*{typed memory; eliminated types}
\For{$t \leftarrow 1$ \KwTo $\budget$}{
  $p_t \sim \prop(\cdot \mid \elim)$ \label{alg:steer}
    \tcp*{steer: avoid eliminated types, \cref{eq:steer}}
  \If{$\mathrm{Guard}_\mem(p_t)$ \label{alg:guard}}{
    \textbf{continue}
      \tcp*{a stored counterexample already refutes $p_t$, \cref{eq:guard}}
  }
  $r \leftarrow \Ok(p_t)$ \tcp*{the expensive verification round}
  \If{$r = \textsf{accept}$}{
    \KwRet $p_t$ \tcp*{$p_t$ passed every $x \in \Xk$, \cref{thm:sound}}
  }
  $x_t \leftarrow r$; $\tau_t \leftarrow \ttype(p_t, x_t)$\;
  $\mem \leftarrow \mem \cup \{(p_t, x_t, \tau_t)\}$;
  $\elim \leftarrow \elim \cup \{\tau_t\}$
    \tcp*{store $+$ eliminate the class}
}
\KwRet $\bot$\;
\end{algorithm}
\fi
